\clearpage
% An unnumbered dated supplement avoids installing the template's per-appendix
% bibliography sections when this addition has no separate bibliography.
\chapter*{截止日后的磁几何研究增补（2026年10月20日）}
\markboth{磁几何研究增补（2026年10月20日）}{磁几何研究增补（2026年10月20日）}
\addcontentsline{toc}{chapter}{截止日后的磁几何研究增补（2026年10月20日）}
\renewcommand{\thetable}{增补\arabic{table}}
\setcounter{table}{0}

本增补为2026年10月20日形成的后续研究更新。报告前文继续以2026年10月13日为截止日，已有实物、仿真、M2.5与完整任务状态均按原证据保留；本增补不表示这些新增内容已在原截止日前完成，也不改变冻结提交版。增补只涉及独立离线磁几何与输出条件，不改生产控制或现场执行许可。

\section*{一、可观测：区分相对磁几何与绝对海床基准}

主动横切将场强序列变成具有空间坐标的磁剖面。半高宽消除整体幅值尺度，但宽度与真实距离的对应仍依赖场形和标定。以参考探头到缆轴的垂距为$z_*$、探头到同一局部海床的高度为$h_*$，轴线埋深满足$d_{\mathrm{axis}}=z_*-h_*$；电缆顶部覆盖厚度还须减去已知外半径。探头高度需要姿态、杆臂与时延改正，不能直接用压力水深或艇体高度代替。

补充研究在已知平行三芯、有限响应与校准相对轨迹下，分别计算横向坐标$c$与$z_*$的有限剖面区间，再接入独立海床参考。共同绝对高度偏差可以不改变相对磁几何，因而重复同基准磁扫不能自行识别其真实埋深。相反，扫描内漂移或姿态相关杆臂错误会改变相对轨迹，不能靠变量换名消除。

\section*{二、可置信：分项判定与参考生命周期}

求解先检查声明模型的绝对残差，再检查各参数允许解释的分歧、先验边界、数值分辨率与空间支撑。有限网格与多初值不保证穷尽全局解，单参数区间也不是无条件联合95\%保证。采用已知回流辅助分解时，回流几何同样是合成先验；该条件不能从电缆内芯资料自动推出。

独立高度参考只在同原点、同缆段、时效和标定成立时形成绝对输出；相关误差未建模时不默认独立。表\ref{tab:middle-round2-dependency}归纳实际实现的更新规则。这里的独立组合不引入额外轴深先验，也不让参考反向约束磁参数，因此不等价于任意联合后验融合。

\begin{table}[htbp]
\centering
\caption{独立磁几何与海床参考的分项输出规则}\label{tab:middle-round2-dependency}
\begin{tabularx}{\linewidth}{@{}p{4.3cm}X@{}}
\toprule
当前证据变化 & 输出处理\\
\midrule
磁几何充分，海床参考缺失 & 合格的$c,z_*$可保留，轴深和覆盖厚度为空\\
独立参考接入或单独更新 & 利用相对区间传播，再核验绝对输出容差\\
参考到期或撤回 & 撤销绝对输出，不把旧值当作当前深度\\
相对轨迹、模型、共同标定或路段改变 & 旧磁集合失效，重新配准、求解或采样\\
未声明轨迹或独立参考偏差 & 仍可能误发布，需真实标定及独立核验\\
\bottomrule
\end{tabularx}
\end{table}

合理工程基线允许相同剖面、分项、缓存与参考传播，因此本轮将这些关系作为可实施的条件管理，不宣称已形成独特精度优势或已成立的专利创造性。完整开发扫描可以不补扫；单侧空间支撑不足时，只有取得真正新观测并正确配准后才重新判断，无安全测线或预算时保持无效。

\section*{三、可行动：保持原主线与后续验收顺序}

新增离线原型未接入默认ROS巡检工厂，不替代原70 m核心段仿真或真实TMR、Jetson及PC104子链证据。隐蔽参考偏差和相对轨迹错误仍可产生小残差、窄区间却错误的数值；质量检查通过不能直接开放埋深或实艇执行。技术交底目前为内部草稿，申请主体、发明人贡献、创造性和提交时机另行确认。

下一阶段固定现有主动观测与分项许可基线，先在经独立核实的短直目标段验证横向巡检；需要埋深时，补充可溯源横切、响应标定、探头外参与同步海床参考，再保留独立最终测试数据。目标平台同次输入来源、实际下发、板端接收及撤回后运动继续按原工作包验收。零命令不表示艇已停住，新增离线结果不改变M2.5和完整水下任务尚未完成的结论。
